%2multibyte Version: 5.50.0.2890 CodePage: 1254
% \titlecontents{chapter}[0pt]{\addvspace{1pc}\itshape}{\contentsmargin{0pt}
% \bfseries\makebox[0pt][r]{\large\thecontentslabel\enspace}\large}
% {\contentsmargin{0pt}\large}{\hfill\contentspage}[\addvspace{.5pc}]
%\setlength{\oddsidemargin}{15mm}
%\setlength{\evensidemargin}{-5mm} 
%\theoremheaderfont{\bfseries}
%   \theorembodyfont{\rmfamily}
%   \theorembodyfont{\sffamily}


\documentclass[12pt,a4paper,turkish]{book}
%%%%%%%%%%%%%%%%%%%%%%%%%%%%%%%%%%%%%%%%%%%%%%%%%%%%%%%%%%%%%%%%%%%%%%%%%%%%%%%%%%%%%%%%%%%%%%%%%%%%%%%%%%%%%%%%%%%%%%%%%%%%%%%%%%%%%%%%%%%%%%%%%%%%%%%%%%%%%%%%%%%%%%%%%%%%%%%%%%%%%%%%%%%%%%%%%%%%%%%%%%%%%%%%%%%%%%%%%%%%%%%%%%%%%%%%%%%%%%%%%%%%%%%%%%%%
\usepackage{amsfonts}
\usepackage{aufbemaster}
\usepackage{ulem}
\usepackage[nodotinlabels]{titletoc}
\usepackage{theorem}
\usepackage{verbatim}
\usepackage{tikz}
\usepackage[T1]{fontenc}
\usepackage[utf8]{inputenc} % also latin5
\usepackage[turkish]{babel}

%TCIDATA{OutputFilter=LATEX.DLL}
%TCIDATA{Version=5.50.0.2890}
%TCIDATA{Codepage=1254}
%TCIDATA{<META NAME="SaveForMode" CONTENT="1">}
%TCIDATA{BibliographyScheme=Manual}
%TCIDATA{Created=Tuesday, February 03, 2015 18:33:51}
%TCIDATA{LastRevised=Wednesday, February 04, 2015 14:15:16}
%TCIDATA{<META NAME="GraphicsSave" CONTENT="32">}
%TCIDATA{<META NAME="DocumentShell" CONTENT="Theses\SW\Ankara University Thesis">}
%TCIDATA{CSTFile=Ankara Universitesi FBE Master Tezi.cst}

\theorembodyfont{\sffamily}
\newtheorem{theorem}{Teorem}[chapter]
\newtheorem{acknowledgement}{Te\c{s}ekkür}
\newtheorem{algorithm}{Algoritma}[chapter]
\newtheorem{axiom}{Aksiyom}[chapter]
\newtheorem{case}{Durum}[chapter]
\newtheorem{claim}{\.{I}ddia}[chapter]
\newtheorem{conclusion}{Yarg\i}[chapter]
\newtheorem{condition}{Ko\c{s}ul}[chapter]
\newtheorem{conjecture}{Kestirim}[chapter]
\newtheorem{corollary}{Sonu\c{c}}[chapter]
\newtheorem{criterion}{Kriter}[chapter]
\newtheorem{definition}{Tan\i m}[chapter]
\newtheorem{example}{\"{O}rnek}[chapter]
\newtheorem{exercise}{Al\i \c{s}t\i rma}[chapter]
\newtheorem{lemma}{Lemma}[chapter]
\newtheorem{notation}{Notasyon}[chapter]
\newtheorem{problem}{Problem}[chapter]
\newtheorem{proposition}{\"{O}nerme}[chapter]
\newtheorem{remark}{Uyar\i}[chapter]
\newtheorem{solution}{\C{c}\"{o}züm}[chapter]
\newtheorem{summary}{\"{O}zet}[chapter]
\newenvironment{proof}[1][\.{I}spat]{\noindent \textbf{#1.} }{\  \rule{0.5em}{0.5em}}
\titlecontents{chapter}[20pt]{\addvspace{0pc}\bfseries}{\contentsmargin{0pt}
\bfseries \makebox[0pt][r]{\bfseries \thecontentslabel \enspace}\bfseries}
{\contentsmargin{0pt}\bfseries}{\hfill \contentspage}[\addvspace{.1pc}]
\dottedcontents{chapter}[2.0em]{\bfseries}{1.8em}{1.2pc}
\titlecontents{section}[20pt]{\addvspace{0pc}\bfseries}{\contentsmargin{0pt}
\bfseries \makebox[0pt][r]{\bfseries \thecontentslabel \enspace}\bfseries}
{\contentsmargin{0pt}\bfseries}{\hfill \contentspage}[\addvspace{.1pc}]
\dottedcontents{section}[2.0em]{\bfseries}{1.8em}{0.6pc}
\titlecontents{subsection}[20pt]{\addvspace{0pc}\bfseries}{\contentsmargin{0pt}
\bfseries \makebox[0pt][r]{\bfseries \thecontentslabel \enspace}\bfseries}
{\contentsmargin{0pt}\bfseries}{\hfill \contentspage}[\addvspace{.1pc}]
\dottedcontents{subsection}[3.1em]{\bfseries}{2.8em}{0.6pc}
\input{tcilatex}

\begin{document}


\begin{tezisim}
TEZİNİZİN ADINI BURAYA YAZINIZ (TAMAMI BÜYÜK HARF)
\end{tezisim}

\begin{tezadonay}
Onay Sayfasına Girecek Tez Adını Buraya Yazınız (Baş Harfler Büyük)
\end{tezadonay}

\begin{thesisname}
THE TITLE OF THE THESIS
\end{thesisname}

\begin{bolumad}
Matematik
\end{bolumad}

\begin{bolumadb}
MATEMATIK
\end{bolumadb}

\begin{department}
Mathematics
\end{department}

\begin{derecead}
Doktora
\end{derecead}

\begin{dereceadb}
DOKTORA/YÜKSEK LİSANS
\end{dereceadb}

\begin{degreename}
Ph.D./Master
\end{degreename}

\begin{danisman}
 Prof. Dr. (Doç.Dr. / Dr. Öğr.Üyesi) Nuri ÖZALP
\end{danisman}

\begin{supervisor}
Prof.Dr. (Assoc.Prof.Dr. /Asst. Prof.Dr.) Nuri ÖZALP
\end{supervisor}

\begin{yazar}
Adınız SOYADINIZ
\end{yazar}

\begin{ay}
Ocak
\end{ay}

\begin{aying}
January
\end{aying}

\begin{yil}
2008
\end{yil}

\begin{dogumtarihi}
Gün/Ay/Yıl (01/01/1990 gibi)
\end{dogumtarihi}

\begin{dogumyeri}
...
\end{dogumyeri}

\begin{medenidurumu}
Bekar/Evli
\end{medenidurumu}

\begin{yabancidil}
Yabancı dil (veya diller)
\end{yabancidil}

\begin{lise}
... Lisesi
\end{lise}

\begin{lisans}
... Üniversitesi, ... Bölümü
\end{lisans}

\begin{yl}
...Üniversitesi (Yoksa bu satırı siliniz)
\end{yl}

\begin{kurum}
...(Yoksa bu satırı siliniz )
\end{kurum}

\begin{yayinlar}
...(Yoksa bu satırı siliniz)
\end{yayinlar}

\begin{tarih}
Gün/Ay/Yıl
\end{tarih}

\begin{oybircok}
oy birliği(veya oy çokluğu)
\end{oybircok}

\begin{mudur}
Prof. Dr. Adı SOYADI
\end{mudur}

\begin{juria}
Jüri Üyesi Ünvanı Adı SOYADI

\hspace{2cm}Ankara Üniversitesi, Matematik Anabilim Dalı 
\end{juria}

\begin{jurib}
Jüri Üyesi Ünvanı Adı SOYADI

\hspace{2cm}... Üniversitesi, ... Anabilim Dalı 
\end{jurib}

\begin{juric}
Jüri Üyesi Ünvanı Adı SOYADI

\hspace{2cm}... Üniversitesi, ... Anabilim Dalı 
\end{juric}

\begin{jurid}
Jüri Üyesi Ünvanı Adı SOYADI

\hspace{2cm}... Üniversitesi, ... Anabilim Dalı 
\end{jurid}

\begin{jurie}
Jüri Üyesi Ünvanı Adı SOYADI

\hspace{2cm}... Üniversitesi, ... Anabilim Dalı 
\end{jurie}

\begin{tezozet}
{\small Buraya tez özeti yazılacaktır. En fazla 250 kelime ile
tek sayfaya sığacak şekilde yazınız. Sığdırmada sıkıntı yaşarsanız aufbemaster.sty dosyasında ÖZET SAYFASI kısmındaki \TEXTsymbol{\backslash}vskip 0.4in değerlerini daha küçük alarak (örneğin 0.3in ) araları sıklaştırabilirsiniz. }
\end{tezozet}

\begin{anahtarkelime}
{\small Anahtar kelimeler buraya yazılacak.}
\end{anahtarkelime}

\begin{tezozeting}
{\small Write your abstract here.. }
\end{tezozeting}

\begin{keyword}
{\small Anahtar kelimelerin \.{I}ngilizcesi buraya yaz\i lacak. }
\end{keyword}

\begin{tsk}
Buradan başlayarak teşekkürünüzü yazınız. Tarih ve isminiz otomatik olarak eklenecektir. 

\end{tsk}

\begin{sayfa}
\pageref{son}
\end{sayfa}

\begin{kapaksayfa}
\end{kapaksayfa}

\onaydok %yüksek lisans için \onayyl yi aktif yapınız
%\onayyl
\begin{etik}
\end{etik}

\begin{ozet}
\end{ozet}

\begin{abstractpage}
\end{abstractpage}

\begin{tesekkur}
\end{tesekkur}

\begin{thetableofcontents}
\end{thetableofcontents}

\begin{nomenclature}
\end{nomenclature}

%\newpage
%\thispagestyle{plain}
%\begin{singlespace}

\noindent 
\begin{tabular}{ll}
$\sum $ & Toplam \\ 
$\Gamma (x)$ & Gamma fonksiyonu \\ 
$H_{2}O$ & Su \\ 
SP \  \  \  \  \  \  \  \  \  \  \  \  \  \  \  \  \  \  \  \  & Scientific Programming \\ 
LGN & Lateral Geniculate Nucleus \\ 
ODC & Ocular Dominance Column \\ 
ODF & Olasılık dağılım fonksiyonu \\ 
$\mathbb{R}$ & Reel uzay \\ 
\end{tabular}

% \end{singlespace}
\newpage

\begin{thelistoffigures}
\end{thelistoffigures}

\begin{thelistoftables}
\end{thelistoftables}

\newpage

%Arkalı-önlü çıktı için}{\setlength{\oddsidemargin}{15mm}} }%
\setlength{\oddsidemargin}{15mm}
%Bu iki  koda dokunmayınız{\setlength{\evensidemargin}{-5mm}}}%
\setlength{\evensidemargin}{-5mm}%


\chapter{G\.{I}R\.{I}\c{S}}

\noindent Bu  template dosyası Ankara Üniversitesi Fen Bilimleri Enstitüsü Tez formatına göre hazırlanmış olup serbestçe kullanılabilir. Nuri ÖZALP \\

Buradan başlayarak tezinizi yazabilirsiniz





\chapter{ENST\.{I}T\"{U} FORMATINA G\"{O}RE TEZ YAZIMI}



\bigskip

Bir bilgisayar programı yazılımı gibi durumlarda metnin orjinali
ile aynen yazılmasını istiyorsanız \textrm{Verbatim} tagı
kullanınız
\begin{verbatim}
 
#include <iostream>        // < > standart kutuphaneler icin.
void main(void)            // 
{
  cout << "Merhaba dunya.";  // Cikti }
 
 
\end{verbatim}


Teoremler Kesimlere göre numaralanmalıdır

\begin{theorem}
Bu bir teoremdir.
\end{theorem}

\begin{theorem}[Ana teorem]
Bir başka teorem. (Etiketli)
\end{theorem}

\begin{proof}
Bu teoremin ispatı için \cite{ozalp} ya bakınız. Dikkat edilirse
burada bir atıf yapılmıştır. 
\end{proof}

\begin{lemma}
Bu bir lemmad\i r ve ispat\i \ i\c{c}in \cite{ozalp2} ye bak\i n\i z.
\end{lemma}

\begin{corollary}
Bu bir sonuçtur.
\end{corollary}

\begin{proposition}
Bu bir Önermedir.
\end{proposition}

\begin{corollary}
Bu da bir sonuçtur.
\end{corollary}

\begin{remark}
Bu bir uyarıdır \  \ 
\end{remark}

\chapter{D\.{I}\u{G}ER B\"{O}L\"{U}MLER}

Sonraki bölümlemleriniz

\chapter{TARTI\c{S}MA VE SONU\c{C}}

Son bölüm TARTIŞMA  ve SONUÇ adı altında olmalıdır

Kaynaklar alfabetik sıraya göre aşağıdaki formatta olmalıdır

\bigskip

\begin{thebibliography}{(Demirci ve Ozalp 2012)}
\bibitem[(Demirci ve Ozalp 2012)]{d1} Demirci, E., Ozalp, N., 2012. \textsl{%
A Method for Solving Differential Equations of Fractional Order}, Journal of
Computational and Applied Mathematics, 236, 2754--2762.

\bibitem[(Hirshfelder vd. 1964)]{r4} Hirshfelder, J. O., Curtis, C. F. and
Bird, R. B\textbf{.,} 1964. \textsl{Molecular Theory of \ Gases and Liquids}%
. John Wiley and Sons, Inc.

\bibitem[(\"{O}zalp 2006)]{ozalp} \"{O}zalp, N., 2006. \textsl{Matematiksel
Modelleme}\textit{.} Gazi Yay\i nlar\i , Ankara.

\bibitem[(\"{O}zalp 2007)]{ozalp2} \"{O}zalp, N., 2007. \textsl{Tez Yaz\i m
Klavuzu}, Pdf dosyas\i , A\"{U}FBE.
\end{thebibliography}

\bigskip

\chapter{\.{I}LK EK ADI}

Tex alan\i nda \TEXTsymbol{\backslash}appendix (ek) fragmenti sadece bir kez
kullan\i l\i r. Ard\i \c{s}\i k ekler B\"{o}l\"{u}m d\"{u}\u{g}mesi se\c{c}%
imi ile verilir. EK b\"{o}l\"{u}mleriniz yok ise bu kesimleri siliniz.

\ 

\chapter{\.{I}K\.{I}NC\.{I} EK ADI}

Ek kesimleriniz yoksa bu b\"{o}l\"{u}mleri siliniz

\bigskip

\begin{ozgecmis}
\label{son}
\end{ozgecmis}

\end{document}
